% 本书术语表（附录 B），由 work/glossary.tsv 经 work/tools/build_glossary.py 生成，勿手改。
% 数据源：work/glossary.tsv（冻结表）；生成脚本：work/tools/build_glossary.py。
% 注意：日文列的中点 ・（U+30FB）在 SimSun 中缺字，已统一改为 ·（U+00B7）；
%       日文假名、长音符 ー、小假名 ャ ュ ョ ッ 在 SimSun 中可正常显示（已实测）。
%       数学符号请用 $...$，正文勿直接写 ≈、≤ 等（Latin Modern 缺字）。

\chapter{术语表（日汉对照）}

\noindent
本表收录全书术语。日文列保留原书写法，OCR 变体以斜线分隔；人名一律保留原文，不音译。
说明列给出首次出现处的处理、译法依据或歧义提示。

\small
\renewcommand{\arraystretch}{1.15}
\begin{longtable}{@{}>{\raggedright\arraybackslash}p{0.29\linewidth}>{\raggedright\arraybackslash}p{0.26\linewidth}>{\raggedright\arraybackslash}p{0.38\linewidth}@{}}
\toprule
\textbf{日文} & \textbf{中文} & \textbf{说明} \\
\midrule
\endfirsthead
\toprule
\textbf{日文} & \textbf{中文} & \textbf{说明} \\
\midrule
\endhead
\bottomrule
\endfoot
\multicolumn{3}{@{}l}{\textbf{一、书名、主题与核心概念}}\\
\midrule
ミラー対称性 & \mbox{镜像对称} & 全书主题；不音译为“米勒对称性” \\
古典的ミラー対称性 & \mbox{经典镜像对称} & classical mirror symmetry \\
位相的ミラー対称性 & \mbox{拓扑镜像对称} & topological mirror symmetry \\
同調ミラー対称性 / ホモロジー的ミラー対称性 / ホモロジカル·ミラー対称性 & \mbox{同调镜像对称} & homological mirror symmetry，缩写 HMS \\
ミラー対称性予想 & \mbox{镜像对称猜想} & mirror symmetry conjecture \\
ミラー対（ミラーペア） & \mbox{镜像对} & mirror pair \\
ミラー多様体 & \mbox{镜像流形} & mirror manifold \\
ミラー変換 & \mbox{镜像变换} & mirror map \\
ミラー写像 & \mbox{镜像映射} & mirror map（第 6 章原书用此词；余章作「镜像变换」） \\
カラビ\textendash{}ヤウ多様体 / カラピーヤウ多様体 / カラビヤウ多様体 / カラピヤウ & Calabi--Yau \mbox{流形} & 首次出现括注原文 \\
カラビ\textendash{}ヤウ条件 / カラピーヤウ条件 & Calabi--Yau \mbox{条件} &  \\
超弦理論 & \mbox{超弦理论} &  \\
コンパクト化 & \mbox{紧化} & compactification \\
内部空間 & \mbox{内部空间} & internal space \\
グロモフ\textendash{}ウィッテン不変量 / グロモフーウィッテン & Gromov--Witten \mbox{不变量} & G--W 不变量 \\
フレアーホモロジー & Floer \mbox{同调} &  \\
フレアー理論 & Floer \mbox{理论} &  \\
サイバーグ·ウィッテン理論 & Seiberg--Witten \mbox{理论} & （第 10 章 p.142；tex 写 Seiberg--Witten） \\
ヘガードフレアー理論 & Heegaard Floer \mbox{理论} & （第 10 章 p.142） \\
D-加群 / D 模 & \mbox{D-模} & D-module \\
ホモロジー代数 & \mbox{同调代数} &  \\
共形場の理論 & \mbox{共形场论} & CFT \\
位相的場の理論 & \mbox{拓扑场论} & TFT \\
量子コホモロジー環 & \mbox{量子上同调环} & quantum cohomology ring（第 2 章起） \\
量子コホモロジー & \mbox{量子上同调} & quantum cohomology（第 8 章 p.116） \\
積構造 & \mbox{积结构} & product structure（第 8 章 p.116） \\
量子変形 & \mbox{量子形变} & quantum deformation（第 8 章 p.116） \\
フーリエ級数 & Fourier \mbox{级数} & Fourier series（第 8 章 p.112） \\
2次形式 & \mbox{二次形式} & quadratic form（第 8 章 p.113） \\
数え上げ幾何学 / 数え上げ幾何 & \mbox{计数几何} & enumerative geometry（第 2 章） \\
位相的弦理論 & \mbox{拓扑弦理论} &  \\
位相的な弦 / 位相的弦 & \mbox{拓扑弦} & topological string（第 10 章 p.131） \\
位相的シグマ模型 & \mbox{拓扑} sigma \mbox{模型} &  \\
非線形シグマ模型 & \mbox{非线性} sigma \mbox{模型} &  \\
トーリック幾何 & \mbox{环簇几何} & toric geometry \\
トーリック多様体 & \mbox{环簇} & toric variety \\
トーリックファノ多様体 & \mbox{环簇} Fano \mbox{流形} & toric Fano variety（第 7 章） \\
代数的トーラス & \mbox{代数环面} & algebraic torus（第 7 章） \\
実トーラス & \mbox{实环面} & real torus（第 7 章） \\
半次元のトーラス & \mbox{半维环面} & half-dimensional torus（第 7 章，SYZ 纤维） \\
完全交叉 & \mbox{完全交} & complete intersection（第 7 章） \\
シンプレクティック商 & \mbox{辛商} & symplectic quotient（第 7 章） \\
反射的多面体 & \mbox{反射多面体} & reflexive polytope \\
SYZ猜想 & SYZ \mbox{猜想} & 原书用 Strominger--Yau--Zaslow 表述 \\
\multicolumn{3}{@{}l}{\textbf{二、基础数学用语}}\\
\midrule
多様体 & \mbox{流形} & manifold \\
写像 & \mbox{映射} & map \\
関数 & \mbox{函数} & function \\
空間 & \mbox{空间} &  \\
集合 & \mbox{集合} &  \\
群 & \mbox{群} &  \\
群化 & \mbox{群化} & groupification（第 10 章 p.141） \\
環 & \mbox{环} & ring \\
体 & \mbox{域} & field \\
加群 & \mbox{模} & module \\
拡大群 & Ext \mbox{群} & 原书「0 次以外的扩张群」译为「除 0 次以外的 Ext 群」（第 11 章） \\
極大コーエン\textendash{}マコーレー加群 & \mbox{极大} Cohen--Macaulay \mbox{模} & maximal Cohen--Macaulay module（第 11 章） \\
行列因子化 & \mbox{矩阵因子化} & matrix factorization（第 11 章） \\
層 & \mbox{层} & sheaf \\
束 & \mbox{丛} & bundle \\
圏 & \mbox{范畴} & category \\
対象 & \mbox{对象} & object（范畴论；第 10 章 p.127） \\
コンパクトな対象 & \mbox{紧对象} & compact object（第 11 章） \\
余表現可能関手 & \mbox{余可表函子} & co-representable functor（第 11 章） \\
セール関手 & Serre \mbox{函子} & Serre functor（第 11 章） \\
擬同型 & \mbox{拟同构} & quasi-isomorphism（第 10 章 p.127） \\
弱同値 & \mbox{弱等价} & weak equivalence（第 10 章 p.127） \\
函手 & \mbox{函子} & functor \\
複体 & \mbox{复形} & complex \\
鎖複体 & \mbox{链复形} & chain complex \\
完全列 & \mbox{正合列} & exact sequence \\
準同型 / 準同型写像 & \mbox{同态} / \mbox{同态映射} & homomorphism \\
同型 & \mbox{同构} & isomorphism \\
核 & \mbox{核} & kernel \\
像 & \mbox{像} & image \\
商 & \mbox{商} & quotient \\
直和 / 直和分解 & \mbox{直和} / \mbox{直和分解} &  \\
テンソル積 & \mbox{张量积} & tensor product \\
基底 & \mbox{基} & basis \\
次元 & \mbox{维} & dimension \\
階数 & \mbox{秩} & rank \\
固有値 & \mbox{特征值} & eigenvalue \\
固有関数 & \mbox{特征函数} & eigenfunction（第 5 章） \\
内積 & \mbox{内积} & inner product \\
計量 & \mbox{度量} & metric \\
位相 & \mbox{拓扑} & topology \\
位相（phase） & \mbox{相位} & 物理语境：位相空间=相空间、位相因子=相位因子；与 topology 义分列 \\
代数的位相幾何学 / 位相幾何学 & \mbox{代数拓扑} & algebraic topology（第 10 章 p.143） \\
トポロジー & \mbox{拓扑} & topology \\
位相空間 & \mbox{拓扑空间} & 数学义（topological space）；与下条「相空间」区分 \\
相空間 & \mbox{相空间} & phase space（第 8 章哈密顿形式） \\
位相不変量 & \mbox{拓扑不变量} &  \\
ベクトル & \mbox{向量} & vector；OCR 常作 ベクトノレ \\
ベクトル空間 & \mbox{向量空间} &  \\
ベクトル束 & \mbox{向量丛} &  \\
直線束 & \mbox{线丛} & line bundle \\
コンパクト & \mbox{紧} & compact \\
パラメータ & \mbox{参数} & parameter \\
モデル & \mbox{模型} & model；Aモデル／Bモデル = A 模型／B 模型 \\
クラス & \mbox{类} & class \\
タイプ & \mbox{类型} & type \\
データ & \mbox{数据} & data \\
ゼロ & \mbox{零} & zero \\
イデアル & \mbox{理想} & ideal \\
ポテンシャル & \mbox{势} & potential \\
プレポテンシャル & \mbox{预势} & prepotential（第 1 章） \\
レベル（WZW 等） & \mbox{水平} & level；Chern--Simons 的 k 亦称水平（p.154） \\
レベル（一般） & \mbox{层面} & 「…のレベルでは」→「在…层面上」 \\
ノード & \mbox{结点} & node；A1 型奇点 xy−zw=0；第 4 章正文统一作「结点」 \\
ループ / ループ空間 & \mbox{圈} / \mbox{环路空间} & loop \\
T 双対 / T-双対 & T \mbox{对偶} & T-duality（第 12 章） \\
1-サイクル & \mbox{1-闭链} & 1-cycle（第 9 章 p.118；OCR 常作「1-サイクル」） \\
モジュラー関数 & \mbox{模函数} & modular function（第 9 章 p.119） \\
模空間 & \mbox{模空间} & moduli space（通用条；第 4 章\textendash{}第 12 章共 38 次） \\
旗流形 & \mbox{旗流形} & flag manifold（第 9 章 p.119） \\
ミラーパートナー & \mbox{镜像伙伴} & mirror partner（第 9 章 p.119） \\
記号的分解 & \mbox{符号分解} & symbolic decomposition（Frobenius 符号分解） \\
基本解行列 & \mbox{基本解矩阵} & fundamental solution matrix（第 9 章 p.112） \\
部分 D-加群 & \mbox{子} \mbox{D-模} & D-submodule（原书「部分 D-加群」，第 9 章 p.118） \\
半無限次元 & \mbox{半无穷维} & half-infinite dimensional（第 8 章 p.98 起） \\
層係数コホモロジー & \mbox{层系数上同调} & cohomology with coefficients in a sheaf \\
ディラックの海 & \mbox{狄拉克之海} & Dirac sea（第 8 章 p.98） \\
力学系 & \mbox{动力系统} & dynamical system（第 9 章 p.113） \\
正準方程式 & \mbox{正则方程} & canonical equation（第 8 章 p.105） \\
周期解 & \mbox{周期解} & periodic solution（第 8 章 p.105） \\
佐藤グラスマン多様体 & \mbox{佐藤} Grassmann \mbox{流形} & Sato Grassmannian（第 8 章 p.99） \\
生成元 & \mbox{生成元} & generator（全书 13 次） \\
保型形式 & \mbox{自守形式} & automorphic form；书目例外 第 5 章 p.236 \\
正則直線束 & \mbox{全纯线丛} & holomorphic line bundle（第 1 章起） \\
正則切断 & \mbox{全纯截面} & holomorphic section（全书 5 次） \\
複素化ケーラー·パラメータ / 複化 Kähler パラメータ & \mbox{复化} Kähler \mbox{参数} & complexified Kähler parameter（第 10 章） \\
閉弦 / 開弦 & \mbox{闭弦} / \mbox{开弦} & closed string / open string（第 8 章 p.108） \\
グラスマン多様体 / グラスマニアン & Grassmann \mbox{流形} & Grassmann manifold（佐藤流形，第 8 章） \\
サイクル & \mbox{闭链} & cycle；3サイクル=3-闭链；不用「圈」（圈只用于 loop） \\
消滅サイクル & \mbox{消没闭链} & vanishing cycle \\
カップ / カップ積 & \mbox{杯} / \mbox{杯积} & cup product \\
ファイバー & \mbox{纤维} & fiber \\
ファイブレーション & \mbox{纤维化} & fibration \\
キャンセル & \mbox{消去} & cancellation（临界点消去） \\
スピン & \mbox{自旋} & spin（粒子自旋） \\
スピン構造 & spin \mbox{结构} & spin structure \\
スピノル & \mbox{旋量} & spinor \\
インスタントン & \mbox{瞬子} & instanton \\
エルミート計量 & Hermite \mbox{度量} & Hermitian metric \\
ユークリッド空間 & Euclid \mbox{空间} &  \\
ローレンツ内積 & Lorentz \mbox{内积} &  \\
ガウス分解 & Gauss \mbox{分解} &  \\
コーシー\textendash{}リーマン方程式 & Cauchy--Riemann \mbox{方程} &  \\
ド·ラームコホモロジー & de Rham \mbox{上同调} &  \\
ラプラス方程式 / ラプラシアン & Laplace \mbox{方程} / Laplace \mbox{算子} &  \\
ポアソン構造 & Poisson \mbox{结构} &  \\
モーメント写像 & \mbox{矩映射} & moment map \\
トポス & \mbox{拓扑斯} & topos（第 10 章） \\
ディリクレ & Dirichlet &  \\
アナロジー & \mbox{类比} & analogy \\
ソリトン / ソリトン方程式 & \mbox{孤子} / \mbox{孤子方程} & soliton \\
ヤング図 & Young \mbox{图} & Young diagram；第 1 章曾作「杨图」，v3 统一为 Young 图 \\
母関数 & \mbox{母函数} & generating function \\
量子補正 & \mbox{量子修正} & quantum correction \\
繰り込み & \mbox{重整化} & renormalization（第 10 章 p.144） \\
湯川結合 & Yukawa \mbox{耦合} & Yukawa coupling \\
相関関数 & \mbox{关联函数} & correlation function \\
演算子 / 作用素 & \mbox{算子} & operator \\
場の理論 & \mbox{场论} &  \\
状態 & \mbox{态} & state \\
中間状態 & \mbox{中间状态} & intermediate state（第 10 章 p.144） \\
双対性 & \mbox{对偶性} & duality \\
\multicolumn{3}{@{}l}{\textbf{三、流形、复几何与代数几何}}\\
\midrule
シンプレクティック多様体 & \mbox{辛流形} &  \\
シンプレクティック形式 & \mbox{辛形式} &  \\
シンプレクティック構造 & \mbox{辛结构} &  \\
概複素構造 & \mbox{近复结构} & almost complex structure \\
概複素多様体 & \mbox{近复流形} & almost complex manifold（正文作「近复流形」） \\
複素構造 & \mbox{复结构} &  \\
ケーラー多様体 & Kähler \mbox{流形} &  \\
ケーラー計量 & Kähler \mbox{度量} &  \\
ケーラー類 & Kähler \mbox{类} &  \\
ケーラー·パラメータ & Kähler \mbox{参数} &  \\
複素多様体 & \mbox{复流形} &  \\
代数多様体 & \mbox{代数簇} &  \\
射影空間 & \mbox{射影空间} &  \\
複素射影平面 & \mbox{复射影平面} &  \\
重み付き射影空間 & \mbox{加权射影空间} & weighted projective space \\
重み付き & \mbox{带权} & weighted（不可作「带重数」，重数=multiplicity） \\
射の空間 & \mbox{态射空间} & morphism space（A$\infty$ 范畴的 hom 空间） \\
トーラス & \mbox{环面} & torus \\
複素トーラス & \mbox{复环面} &  \\
アーベル多様体 & \mbox{阿贝尔簇} &  \\
代数曲面 & \mbox{代数曲面} &  \\
K3曲面 & K3 \mbox{曲面} &  \\
クインティック & \mbox{五次超曲面} & quintic；五次三维流形 \\
ファノ多様体 & Fano \mbox{流形} &  \\
グラスマニアン & Grassmann \mbox{流形} & Grassmannian \\
クンマー曲面 & Kummer \mbox{曲面} &  \\
エンリケス曲面 & Enriques \mbox{曲面} &  \\
楕円曲線 & \mbox{椭圆曲线} &  \\
楕円ファイブレーション & \mbox{椭圆纤维化} & elliptic fibration \\
ピカール群 & Picard \mbox{群} &  \\
ピカール·フックス方程式 & Picard--Fuchs \mbox{方程} & 第 1 章；OCR 常作 ピカーノレフックス \\
セール双対 & Serre \mbox{对偶} &  \\
リーマン\textendash{}ロッホの定理 & Riemann--Roch \mbox{定理} &  \\
チャーン類 & Chern \mbox{类} &  \\
第一チャーン類 & \mbox{第一} Chern \mbox{类} &  \\
チャーン指標 & Chern \mbox{特征} & Chern character \\
トッド類 & Todd \mbox{类} &  \\
オイラー数 & Euler \mbox{数} &  \\
リッチ曲率 & Ricci \mbox{曲率} &  \\
リッチ平坦 & Ricci \mbox{平坦} &  \\
リッチ·テンソル & Ricci \mbox{张量} &  \\
スカラー曲率 & \mbox{数量曲率} &  \\
アインシュタイン方程式 & Einstein \mbox{方程} &  \\
接続 & \mbox{联络} & connection \\
平坦接続 & \mbox{平坦联络} & flat connection（第 10 章 p.142） \\
反自己双対接続 & \mbox{反自对偶联络} & anti-self-dual connection \\
複素化 & \mbox{复化} & complexification \\
周期 & \mbox{周期} & period；共轭闭链的周期 \\
周期積分 & \mbox{周期积分} & period integral \\
周期写像 & \mbox{周期映射} & period map \\
大域的 & \mbox{整体} & global；作定语作「整体的」 \\
局所パラメータ & \mbox{局部参数} & local parameter \\
複素構造のモジュライ & \mbox{复结构模空间} & moduli of complex structures \\
ケーラー·モジュライ & Kähler \mbox{模空间} & moduli of Kähler structures \\
\multicolumn{3}{@{}l}{\textbf{四、代数几何与极小模型纲领（第 4 章）}}\\
\midrule
極小モデル & \mbox{极小模型} & minimal model \\
極小モデル計画 / 極小モデル·プログラム & \mbox{极小模型纲领} & MMP \\
双有理同値 & \mbox{双有理等价} & birational equivalence \\
同値 & \mbox{等价} & equivalence（第 11 章起统一译「等价」，不译「同值」） \\
Π-安定性 & \mbox{Π-稳定性} & Π-stability（第 11 章；保留希腊字母 Π） \\
小平次元 & Kodaira \mbox{维数} & Kodaira dimension \\
標準束 & \mbox{典范丛} & canonical bundle \\
標準因子 & \mbox{典范除子} & canonical divisor \\
反標準因子 & \mbox{反典范除子} & anticanonical divisor \\
因子（代数幾何） & \mbox{除子} & divisor；标准／反标准／例外／Weil／Cartier 因子均用「除子」 \\
因子（一般） & \mbox{因子} & factor；如 multiple-cover 的修正因子 \\
因子類群 & \mbox{除子类群} & divisor class group \\
例外因子 & \mbox{例外除子} & exceptional divisor \\
ヴェイユ因子 & Weil \mbox{除子} &  \\
カルティエ因子 & Cartier \mbox{除子} &  \\
ネフ（数値的正） & \mbox{nef（数值正）} & nef 无通行中文名，保留原文 \\
豊富 & \mbox{丰富} & ample；丰富性定理中 $K_X$ 半丰沛 \\
豊富性定理 & \mbox{丰富性定理} & abundance theorem（川又\textendash{}宫冈） \\
末端特異点 & \mbox{终端奇点} & terminal singularity \\
標準特異点 & \mbox{典范奇点} & canonical singularity \\
Q分解的 & \mbox{Q-分解} & Q-factorial \\
ブローアップ / プローアップ & \mbox{爆破} & blow-up \\
ブローダウン & \mbox{收缩} & blow-down \\
収縮 & \mbox{收缩} & contraction \\
フリップ & \mbox{翻转} & flip \\
フロップ & flop & 无通行中文名，保留原文 \\
例外曲線 & \mbox{例外曲线} &  \\
有理曲線 & \mbox{有理曲线} &  \\
(−1)曲線 & \mbox{(−1)-曲线} &  \\
交叉数 & \mbox{相交数} & intersection number \\
留数 & \mbox{留数} & residue；Grothendieck 留数照此 \\
環層空間 / 環付き空間 & \mbox{环层空间} & ringed space \\
重複度 & \mbox{重数} & multiplicity；原书作「计入重数」讲，此处从简 \\
種数 & \mbox{亏格} & genus；0 种数 = 0 亏格 \\
単純リー環 & \mbox{单李代数} & simple Lie algebra \\
基本対称式 & \mbox{初等对称多项式} & elementary symmetric polynomial \\
不正則数 & \mbox{非正则数} & irregularity \\
小平埋め込み定理 & \mbox{小平嵌入定理} & Kodaira embedding theorem \\
線織曲面 & \mbox{直纹曲面} & ruled surface \\
超楕円曲面 & \mbox{超椭圆曲面} & hyperelliptic surface \\
指数 1 被覆 & \mbox{指数} 1 \mbox{覆盖} & index 1 cover \\
特異カラビ\textendash{}ヤウ多様体 & \mbox{带奇点的} Calabi--Yau \mbox{流形} & singular Calabi--Yau manifold \\
特異カラビ\textendash{}ヤウ·トランジション & \mbox{带奇点的} Calabi--Yau \mbox{跃迁} & singular CY transition \\
有理楕円曲面 & \mbox{有理椭圆曲面} & rational elliptic surface \\
ファイバー積 & \mbox{纤维积} & fiber product \\
双斉次座標 & \mbox{双齐次坐标} & bihomogeneous coordinates \\
留数形式 & \mbox{留数形式} & residue form \\
小特異点解消 & \mbox{小消解} & small resolution \\
特異点解消 & \mbox{奇点消解} & resolution of singularities \\
クレパント解消 & crepant \mbox{消解} & crepant resolution；crepant 不音译 \\
非可換クレパント解消 & \mbox{非交换} crepant \mbox{解消} &  \\
非可換特異点解消 & \mbox{非交换奇点消解} & 保持标准丛的解消；与 crepant 解消（非交换情形）区分（第 11 章） \\
コニフォールド & \mbox{锥形流形} & conifold \\
コニフォールド·トランジション & \mbox{锥形流形跃迁} & conifold transition \\
幾何学的遷移 & \mbox{几何跃迁} & geometric transition \\
極小特異点解消 & \mbox{极小消解} &  \\
\multicolumn{3}{@{}l}{\textbf{五、奇点、ADE 与环簇（第 3、5、7、11 章）}}\\
\midrule
特異点 & \mbox{奇点} & singularity；原书同句连用时不再重复译 \\
超曲面特異点 & \mbox{超曲面奇点} & hypersurface singularity \\
孤立特異点 & \mbox{孤立奇点} & isolated singularity \\
群作用付き孤立特異点 & \mbox{带群作用的孤立奇点} &  \\
単純特異点 & \mbox{单奇点} & simple singularity，即 ADE 型奇点；与「单李代数」用字一致 \\
ADE型特異点 & ADE \mbox{型奇点} &  \\
商特異点 & \mbox{商奇点} & quotient singularity \\
重み付き斉次超曲面特異点 & \mbox{加权齐次超曲面奇点} &  \\
ミルナー数 & Milnor \mbox{数} & Milnor number \\
ミルナー·モノドロミー & Milnor \mbox{单值性} & Milnor monodromy \\
モノドロミー & \mbox{单值性} & monodromy \\
中間ホモロジー群 & \mbox{中间同调群} & middle homology group \\
ヤコビ環 & Jacobi \mbox{环} & Jacobi ring \\
ヤコビ環のポアンカレ級数 & Jacobi \mbox{环的} Poincaré \mbox{级数} &  \\
ルート系 & \mbox{根系} & root system；OCR 常作 ノレート系 \\
古典ルート系 & \mbox{古典根系} & classical root system（第 11 章） \\
一般ルート系 & \mbox{一般根系} & general root system（第 11 章） \\
ルート格子 & \mbox{根格} & root lattice（第 11 章） \\
交差形式 / 交叉形式 & \mbox{相交形式} & intersection form（第 11 章原书两处印法不一） \\
ワイル群 & Weyl \mbox{群} &  \\
カルタン行列 & Cartan \mbox{矩阵} &  \\
カルタン部分代数 & Cartan \mbox{子代数} &  \\
デインキン図（式） / ディンキン図形 & Dynkin \mbox{图} &  \\
拡大デインキン図形 & \mbox{扩充} Dynkin \mbox{图} & extended（affine）Dynkin diagram \\
経路代数 & \mbox{路径代数} & path algebra \\
ディンキン箙 & Dynkin \mbox{箭图} & quiver（第 3 章 p.42，原书振假名「かご」） \\
拡大ディンキン箙 & \mbox{扩充} Dynkin \mbox{箭图} & affine 型 Dynkin 箭图（第 3 章 p.44） \\
カスプ特異点 & \mbox{尖点奇点} & cusp singularity（第 3 章 p.42） \\
交差形式 & \mbox{相交形式} & intersection form（第 3 章 p.42） \\
オイラー形式 & Euler \mbox{形式} & Euler form（第 3 章 p.42 对照表） \\
グロタンディーク群 & Grothendieck \mbox{群} & （第 3 章 p.42 对照表） \\
三角群 & \mbox{三角群} & triangle group（Schwarz 三角群，第 3 章 p.45） \\
マッカイ対応 / マッケイ対応 & McKay \mbox{对应} &  \\
有限部分群 & \mbox{有限子群} & finite subgroup \\
乗法群 & \mbox{乘法群} & multiplicative group \\
軌道空間 & \mbox{轨道空间} & orbit space \\
軌道体 / オービフォールド & \mbox{轨形} & orbifold \\
軌道体ポアンカレ多項式 & \mbox{轨形} Poincaré \mbox{多项式} &  \\
部分トーラス & \mbox{子环面} & subtorus（第 7 章；书内旧译「部分环面」应统一为「子环面」） \\
アメーバ / アメーパ & \mbox{阿米巴} & amoeba（第 7 章 p.94） \\
反標準束 & \mbox{反典范丛} & anticanonical bundle \\
ホッジ指数 / ホッジ指数の分散 & Hodge \mbox{指数} / Hodge \mbox{指数的方差} & 分散=variance（二阶中心矩）；p141 式 (4) \\
深谷圏 & Fukaya \mbox{范畴} & Fukaya category；深谷为姓氏，保留汉字 \\
原始形式 & \mbox{原始形式} & primitive form（斋藤恭司，第 5 章） \\
半普遍変形 & \mbox{半通用形变} & semi-universal deformation \\
普遍変形 & \mbox{通用形变} &  \\
変形理論 & \mbox{形变理论} & deformation theory \\
平坦構造 & \mbox{平坦结构} & flat structure \\
平坦座標 & \mbox{平坦坐标} & flat coordinates \\
可積分系 & \mbox{可积系统} & integrable system（第 5 章） \\
スケール変換 & \mbox{尺度变换} & scale transformation（第 5 章） \\
フロベニウス構造 & Frobenius \mbox{结构} & Frobenius structure \\
極双対 & \mbox{极对偶} & polar dual \\
アメーパ & \mbox{阿米巴} & amoeba \\
トロピカル幾何 & \mbox{热带几何} & tropical geometry \\
扇 & \mbox{扇} & fan（环簇理论） \\
有理多面錐 & \mbox{有理多面锥} & rational polyhedral cone（第 7 章） \\
双対格子 & \mbox{对偶格} & dual lattice \\
格子 & \mbox{格} & lattice \\
格子偏極 K3 曲面 / 格偏極 K3 曲面 & \mbox{格偏振} K3 \mbox{曲面} & lattice-polarized K3 surface（第 6 章） \\
分数次元 & \mbox{分数维数} & fractional dimension（第 11 章） \\
凸多面体 & \mbox{凸多面体} & convex polytope \\
ニュートン多面体 & \mbox{牛顿多面体} & Newton polytope \\
商トーラス & \mbox{商环面} & quotient torus \\
極大退化 & \mbox{极大退化} & maximal degeneration \\
大半径極限 / ラージ·ラジアス極限 & \mbox{大半径极限} & large radius limit \\
大複素構造極限 & \mbox{大复结构极限} & large complex structure limit \\
トーリック·ファイブレーション & \mbox{环簇纤维化} & toric fibration \\
ミラー対称性と共形場の理論 & \mbox{镜像对称与共形场论} & 第 6 章 \\
\multicolumn{3}{@{}l}{\textbf{六、范畴、同调代数与 D-模（第 9、10、11 章）}}\\
\midrule
連接層 & \mbox{凝聚层} & coherent sheaf \\
準連接層 & \mbox{拟凝聚层} & quasi-coherent sheaf（第 11 章） \\
構造層 & \mbox{结构层} & structure sheaf \\
前層 & \mbox{预层} & presheaf \\
導来圏 & \mbox{导出范畴} & derived category \\
連接層の導来圏 & \mbox{凝聚层的导出范畴} &  \\
三角圏 & \mbox{三角范畴} & triangulated category \\
満充部分三角圏 & \mbox{满三角子范畴} & full triangulated subcategory（第 11 章） \\
ディスティングイッシュト基底 & distinguished basis & 保留英文，首现括注「特别的基」（第 11 章） \\
エクセプショナル·コレクション & exceptional collection & 保留英文（第 11 章） \\
ホモトピー圏 & \mbox{同伦范畴} & homotopy category \\
圏化 / カテゴリー化 & \mbox{范畴化} & categorification \\
A$\infty$圏 / A無限大圏 & A$\infty$ \mbox{范畴} & A-infinity category \\
A$\infty$代数 / A無限大代数 & A$\infty$ \mbox{代数} & A-infinity algebra \\
弱同形 & \mbox{弱同构} & weak isomorphism \\
局所化 & \mbox{局部化} & localization \\
森田同値 & Morita \mbox{等价} & Morita equivalence \\
構造定数 & \mbox{结构常数} & structure constants \\
自然変換 & \mbox{自然变换} & natural transformation \\
圏の同値 & \mbox{范畴等价} & equivalence of categories \\
随伴関手 & \mbox{伴随函子} & adjoint functor \\
導来関手 & \mbox{导出函子} & derived functor \\
写像錐 & \mbox{映射锥} & mapping cone \\
コーン & \mbox{锥} & cone \\
スペクトル系列 & \mbox{谱序列} & spectral sequence \\
非可換幾何学 & \mbox{非交换几何} & noncommutative geometry \\
スキーム & \mbox{概形} & scheme \\
スタック & \mbox{叠} & stack \\
フーリエ\textendash{}向井変換 & \mbox{Fourier\textendash{}向井变换} &  \\
余接束 & \mbox{余切丛} & cotangent bundle \\
接束 & \mbox{切丛} & tangent bundle \\
法束 & \mbox{法丛} & normal bundle \\
余法束 & \mbox{余法丛} & conormal bundle \\
超平面配置 & \mbox{超平面配置} & hyperplane arrangement \\
ホッジ構造 & Hodge \mbox{结构} &  \\
ホッジ数 & Hodge \mbox{数} &  \\
ホッジ対称性 & Hodge \mbox{对称性} &  \\
変動するホッジ構造 / 変化するホッジ構造 & \mbox{变动的} Hodge \mbox{结构} & VHS \\
超幾何微分方程式 & \mbox{超几何微分方程} &  \\
超幾何関数 & \mbox{超几何函数} &  \\
常微分方程式 & \mbox{常微分方程} & ODE \\
確定特異点 & \mbox{正则奇点} & regular singular point \\
不確定特異点 & \mbox{非正则奇点} & irregular singular point \\
微分作用素 & \mbox{微分算子} & differential operator \\
下三角行列 & \mbox{下三角矩阵} &  \\
対角成分 & \mbox{对角元} & diagonal entry \\
量子D-加群 & \mbox{量子} \mbox{D-模} &  \\
ホロノミック & holonomic & 无通行中文名，保留原文 \\
GKZ方程式 & GKZ \mbox{方程} & Gelfand--Kapranov--Zelevinsky 方程（第 9 章） \\
D-プレーン / D-ブレーン & \mbox{D-膜} & D-brane；第 10 章以 D-模／导出范畴解释 \\
プレーン / ブレーン & \mbox{膜} & brane；OCR 常作 プレーン \\
\multicolumn{3}{@{}l}{\textbf{七、辛几何与 Floer 理论（第 8 章）}}\\
\midrule
ラグランジュ部分多様体 / ラグランジアン部分多様体 & Lagrangian \mbox{子流形} & Lagrangian submanifold；第 10 章误作 Lagrange \\
特殊ラグランジアン部分多様体 & \mbox{特殊} Lagrangian \mbox{子流形} & special Lagrangian submanifold \\
特殊ラグランジュ部分多様体 & \mbox{特殊} Lagrangian \mbox{子流形} & special Lagrangian \\
ラグランジュ·ファイブレーション & Lagrangian \mbox{纤维化} & Lagrangian fibration \\
モース関数 & Morse \mbox{函数} &  \\
モース理論 & Morse \mbox{理论} &  \\
勾配ベクトル場 & \mbox{梯度向量场} & gradient vector field \\
勾配 & \mbox{梯度} & gradient \\
積分曲線 & \mbox{积分曲线} & integral curve \\
臨界点 & \mbox{临界点} & critical point \\
ヘシアン & Hessian &  \\
指数 & \mbox{指数} & index \\
境界作用素 & \mbox{边界算子} & boundary operator \\
余核（CoKer / Coker） & \mbox{余核} & cokernel（第 8 章 p.99 原书 CoKer/Coker 混印） \\
キャンセル（同調論） & \mbox{抵消} & 同调论中的抵消；临界点消去仍作「消去」（第 8 章 p.102） \\
ディリクレ境界条件 & Dirichlet \mbox{边界条件} & Dirichlet 边界条件（第 10 章 p.139） \\
総実 & \mbox{全实} & totally real（第 10 章 p.139；「全实子流形」，通行译法） \\
擬正則写像 & \mbox{拟全纯映射} & pseudoholomorphic map \\
正則写像（複素幾何） & \mbox{全纯映射} & holomorphic map \\
正則写像（代数幾何） & \mbox{正则映射} & morphism；与「有理写像→有理映射」配对（tr3 第 4 章反馈） \\
有理写像 & \mbox{有理映射} & rational map \\
調和写像 & \mbox{调和映射} & harmonic map \\
グロモフ\textendash{}ハウスドルフ収束 & Gromov--Hausdorff \mbox{收敛} &  \\
安定写像 & \mbox{稳定映射} & stable map \\
仮想基本類 & \mbox{虚拟基本类} & virtual fundamental class \\
ベアリング & \mbox{配对} & pairing；OCR 常作 ベアリング \\
ハミルトン形式 & Hamilton \mbox{形式} & Hamiltonian formalism \\
ハミルトニアン & Hamilton \mbox{量} & Hamiltonian \\
ラグランジアン & Lagrange \mbox{量} & Lagrangian \\
ラグランジュの汎関数 & Lagrange \mbox{泛函} &  \\
配位空間 & \mbox{位形空间} & configuration space \\
作用 / 作用汎関数 & \mbox{作用量} / \mbox{作用泛函} & action \\
汎関数 & \mbox{泛函} & functional \\
経路積分 & \mbox{路径积分} & Feynman 路径积分 \\
測地線 & \mbox{测地线} & geodesic \\
閉測地線 & \mbox{闭测地线} & closed geodesic \\
変分法 & \mbox{变分法} & calculus of variations \\
変分問題 & \mbox{变分问题} & variational problem \\
極値 & \mbox{极值} & extremum \\
運動方程式 & \mbox{运动方程} & equation of motion \\
同変指数 & \mbox{等变指数} & equivariant index \\
交点数 & \mbox{相交数} & intersection number \\
グロモフ\textendash{}ウィッテン不変量とフレアーホモロジー & Gromov--Witten \mbox{不变量与} Floer \mbox{同调} & 第 8 章主题 \\
\multicolumn{3}{@{}l}{\textbf{八、物理、弦论与共形场论（第 1、2、6、12 章）}}\\
\midrule
弦理論 & \mbox{弦论} &  \\
閉じた弦 & \mbox{闭弦} & closed string \\
開いた弦 & \mbox{开弦} & open string \\
ワールドシート & \mbox{世界面} & worldsheet \\
リーマン面 & \mbox{黎曼面} & Riemann surface \\
共形不変性 & \mbox{共形不变性} & conformal invariance \\
共形次元 & \mbox{共形维数} & conformal dimension \\
中心電荷 & \mbox{中心荷} & central charge \\
超対称性 & \mbox{超对称} & supersymmetry \\
超重力理論 & \mbox{超引力理论} & supergravity \\
ゲージ理論 & \mbox{规范理论} & gauge theory \\
ゲージ粒子 & \mbox{规范粒子} & gauge particle（第 1 章正文用「规范粒子」） \\
ゲージ場 & \mbox{规范场} & gauge field \\
ゲージ変換 & \mbox{规范变换} & gauge transformation \\
量子化 & \mbox{量子化} & quantization \\
分配関数 & \mbox{配分函数} & partition function \\
ブラックホール & \mbox{黑洞} & black hole \\
時空 & \mbox{时空} & spacetime \\
素粒子 & \mbox{基本粒子} & elementary particle \\
統一理論 & \mbox{统一理论} & unified theory \\
量子重力理論 & \mbox{量子引力理论} & quantum gravity \\
重力子 & \mbox{引力子} & graviton（第 1 章；OCR 未检出该词时以「重力子」为准） \\
重力場 & \mbox{引力场} & gravitational field \\
質量 & \mbox{质量} & mass \\
基底状態 & \mbox{基态} & ground state \\
チャージ & \mbox{荷} & charge \\
フェルミオン & \mbox{费米子} & fermion \\
ボゾン / ポソン & \mbox{玻色子} & boson；OCR 常作 ポソン \\
自由フェルミオン場 & \mbox{自由费米子场} & free fermion field \\
シューア関数 & Schur \mbox{函数} &  \\
オペレーター代数 & \mbox{算子代数} & operator algebra \\
頂点作用素 & \mbox{顶点算子} & vertex operator \\
頂点作用素代数 & \mbox{顶点算子代数} & vertex operator algebra \\
ピラソロ代数 / ビラソロ代数 & Virasoro \mbox{代数} &  \\
ピラソロ作用素 / ビラソロ作用素 & Virasoro \mbox{算子} &  \\
ピラソロ予想 & Virasoro \mbox{猜想} & 第 1 章 \\
カイラル·プライマリー & \mbox{手征原初态} & chiral primary；第 1 章作「手征原初态」，指场时作「手征原初场」 \\
反カイラル·プライマリー & \mbox{反手征原初态} & anti-chiral primary；与「手征原初态」同口径 \\
プライマリー場 & \mbox{原初态} & primary field；正文作「原初态」，强调场时作「原初场」 \\
カイラル環 & \mbox{手征环} & chiral ring \\
セクター & \mbox{扇区（区）} & sector；NS セクター = NS 扇区 \\
ホロノミー & \mbox{和乐群} & holonomy；SU(3) ホロノミー = SU(3) 和乐群 \\
ヘテロティック & \mbox{杂化（heterotic）} & heterotic \\
双対偶 & \mbox{双对偶} & dual pair（第 1 章；与 T 对偶不同源） \\
多重項 & \mbox{多重态} & 向量多重态 / 超多重态（vector / hyper multiplet） \\
鏡映関係 & \mbox{镜像关系} & mirror relation \\
カノニカル束 & \mbox{典范丛} & canonical bundle（= 标准丛） \\
特殊幾何 & \mbox{特殊几何} & special geometry \\
多重被覆 & \mbox{多重覆盖} & multiple cover \\
スペクトラルフロー & \mbox{谱流} & spectral flow（第 6 章） \\
R対称性 & R \mbox{对称性} & R-symmetry \\
超共形場の理論 & \mbox{超共形场论} & superconformal field theory \\
モジュラー不変性 & \mbox{模不变性} & modular invariance；作定语作「模不变」（模不变配分函数） \\
モジュラー形式 & \mbox{模形式} & modular form \\
ナレイン格子 & Narain \mbox{格} & Narain lattice \\
自己双対格子 & \mbox{自对偶格} & self-dual lattice \\
共形場の理論とミラー対称性 & \mbox{共形场论与镜像对称} & 第 6 章 \\
位相的パーテックス / 位相的バーテックス / トポロジカル·バーテックス / 位相的頂点 & \mbox{拓扑顶点} & topological vertex \\
ヒッグス場 & Higgs \mbox{场} & 正文未用 \\
エネルギー & \mbox{能量} & energy \\
運動量 & \mbox{动量} & momentum \\
低エネルギー & \mbox{低能} & low energy \\
分配関数のモジュラー不変性 & \mbox{配分函数的模不变性} & 第 1 章 \\
\multicolumn{3}{@{}l}{\textbf{九、拓扑弦与不变量（第 12 章）}}\\
\midrule
ジョーンズ不変量 & Jones \mbox{不变量} & Jones polynomial \\
ジョーンズ\textendash{}ウィッテン不変量 & Jones--Witten \mbox{不变量} &  \\
チャーン\textendash{}サイモンズ理論 & Chern--Simons \mbox{理论} &  \\
ドナルドソン\textendash{}トーマス不変量 & Donaldson--Thomas \mbox{不变量} &  \\
ドナルドソン不変量 & Donaldson \mbox{不变量} &  \\
壁越え公式 & \mbox{跨墙公式} & wall-crossing formula \\
ファインマン·グラフ & Feynman \mbox{图} & Feynman graph \\
構造群 & \mbox{结构群} & structure group \\
ゴパクマール\textendash{}ヴァファ & Gopakumar--Vafa & 第 12 章人名／理论 \\
ポアンカレ双対 & Poincaré \mbox{对偶} &  \\
基本類 & \mbox{基本类} & fundamental class \\
リード·ファンタジー & Reid \mbox{的幻想} & Reid's fantasy（第 4 章） \\
\multicolumn{3}{@{}l}{\textbf{十、论述与措辞}}\\
\midrule
予想 & \mbox{猜想} & conjecture \\
定理 & \mbox{定理} & theorem \\
命題 & \mbox{命题} & proposition \\
補題 & \mbox{引理} & lemma \\
系 & \mbox{推论} & corollary \\
定義 & \mbox{定义} & definition \\
例 & \mbox{例} & example \\
注意 & \mbox{注} & remark \\
仮定 & \mbox{假设} & assumption \\
反例 & \mbox{反例} & counterexample \\
証明 & \mbox{证明} & proof \\
主張 & \mbox{断言} & claim \\
\multicolumn{3}{@{}l}{\textbf{十一、人名对照（西文保留原文，不音译）}}\\
\midrule
グロモフ & Gromov &  \\
ウィッテン & Witten &  \\
フレアー & Floer &  \\
ケーラー & Kähler &  \\
カラビ & Calabi &  \\
ヤウ & Yau & 丘成桐 \\
リーマン & Riemann &  \\
ホッジ & Hodge &  \\
ポアンカレ & Poincaré &  \\
リッチ & Ricci &  \\
オイラー & Euler &  \\
チャーン & Chern &  \\
モース & Morse &  \\
ディラック & Dirac &  \\
フロベニウス & Frobenius &  \\
グロタンディーク & Grothendieck &  \\
ブリースクルン & Brieskorn &  \\
サイデル & Seidel & 人名保留原文（第 3 章） \\
アカンポ & A'Campo & 人名保留原文（第 3 章） \\
シュバルツ & Schwarz & 人名保留原文（第 3 章） \\
アーノルド & Arnold & 人名保留原文（第 3 章） \\
ドゥブロビン & Dubrovin &  \\
アーベル & Abel &  \\
クンマー & Kummer &  \\
ファノ & Fano &  \\
ナレイン & Narain &  \\
シーガル & Segal &  \\
ドナルドソン & Donaldson &  \\
トーマス & Thomas &  \\
キャンデラス & Candelas &  \\
ルアン & Ruan & 人名保留原文（第 2 章/第 7 章） \\
ティアン & Tian & 人名保留原文（第 2 章） \\
シューベルト & Schubert & 人名保留原文（第 2 章） \\
ネクラソフ & Nekrasov &  \\
マクドナルド & Macdonald &  \\
ゲッチェ / ゲッツラー & Göttsche & OCR 两种写法均指同一人 \\
ゲスト & Guest & M. Guest（第 9 章合著者） \\
乙藤隆史 & \mbox{乙藤隆史} & 日本作者名照原 \\
向井 & \mbox{向井} & 向井茂 \\
深谷賢治 & \mbox{深谷賢治} & 日本作者名照原 \\
中島啓 & \mbox{中島啓} & 日本作者名照原 \\
江口徹 & \mbox{江口徹} & 日本作者名照原 \\
マンフォード & Mumford &  \\
ドリーニュ & Deligne &  \\
ヒルツェブルフ & Hirzebruch &  \\
エンリケス & Enriques &  \\
ヴェイユ & Weil &  \\
クライン & Klein & F. クライン \\
カルタン & Cartan &  \\
セール & Serre &  \\
ベズー & Bézout &  \\
アティヤ & Atiyah &  \\
シンガー & Singer &  \\
ヤコビ & Jacobi &  \\
マッカイ & McKay &  \\
リード & Reid & M. Reid \\
シェーン & Schoen & C. Schoen（第 4 章例） \\
ゲップナー & Gepner &  \\
グリーン & Greene & B. Greene \\
プレッサー & Plesser & M. R. Plesser \\
ギベンタリ & Givental &  \\
コンツェヴィッチ & Kontsevich &  \\
ソイベルマン & Soibelman & Y. Soibelman \\
モリソン & Morrison & D. Morrison \\
グロス & Gross & M. Gross \\
バティレフ & Batyrev &  \\
ミルナー & Milnor &  \\
ネーター & Noether & E. Noether \\
レレイ & Leray & J. Leray \\
ホヴァノフ & Khovanov &  \\
川又 & \mbox{川又} & 川又雄二郎 \\
宮岡洋一 & \mbox{宮岡洋一} &  \\
小平邦彦 & \mbox{小平邦彦} & Kodaira \\
佐藤幹夫 & \mbox{佐藤幹夫} & 孤子理论（原书作 soliton theory） \\
齋藤恭司 / 粛藤恭司 & \mbox{齋藤恭司} & 原始形式；OCR 作异体字（见左列变体） \\
小木曽啓示 & \mbox{小木曽啓示} &  \\
岡潔 & \mbox{岡潔} & Oka \\
森田 & \mbox{森田} & 森田等价（Morita） \\
\end{longtable}
\normalsize
